\subsection{User Study}~\label{exp:user_study}
We conducted a Wizard-of-Oz (WoZ) user study to examine whether the proposed state-level query interface is usable with human participants, as a supplementary feasibility check alongside the system-level evaluation. Twelve participants each experienced every domain--condition pair once; we report task success, operation time, NASA Raw-TLX workload, and SAGAT accuracy.

\begin{figure*}[t]
    \centering
    \begin{subfigure}[t]{0.40\textwidth}
        \centering
        \includegraphics[width=\linewidth]{figures/user_study/waste_task_success.png}
        \caption{Waste success.}
        \label{fig:user_waste_task_success}
    \end{subfigure}
    \hspace{0.03\textwidth}
    \begin{subfigure}[t]{0.40\textwidth}
        \centering
        \includegraphics[width=\linewidth]{figures/user_study/tomato_task_success.png}
        \caption{Tomato success.}
        \label{fig:user_tomato_task_success}
    \end{subfigure}
    \vspace{0.5em}

    \begin{subfigure}[t]{0.40\textwidth}
        \centering
        \includegraphics[width=\linewidth]{figures/user_study/waste_operation_time.png}
        \caption{Waste operation time.}
        \label{fig:user_waste_operation_time}
    \end{subfigure}
    \hspace{0.03\textwidth}
    \begin{subfigure}[t]{0.40\textwidth}
        \centering
        \includegraphics[width=\linewidth]{figures/user_study/tomato_operation_time.png}
        \caption{Tomato operation time.}
        \label{fig:user_tomato_operation_time}
    \end{subfigure}
    \caption{Task success and operation time across query strategies. \textbf{No} is omitted from these two measures because they focus on scorable user responses and active user operation. Error bars indicate standard error.}
    \label{fig:user_task_time_results}
\end{figure*}

\subsubsection{Participants}~\label{user:participants}
Twelve participants (7 female and 5 male) took part in the study. Participants ranged in age from 21 to 26 years ($M=23.0$, $SD=1.48$). Each participant received \$20 upon completion of the study.

\subsubsection{Study Design and Conditions}~\label{user:design}
The study employed a within-subject design. Each participant completed ten trials in total, consisting of five waste sorting scenarios and five tomato harvesting scenarios. In each domain, participants experienced five interaction conditions: \textbf{All}, \textbf{No}, \textbf{Ours1}, \textbf{Ours2}, and \textbf{KnowNo}. \textbf{All} requested feedback whenever a query was available, representing continuous human supervision. \textbf{No} performed the task without requesting user feedback. \textbf{Ours1} used the proposed uncertainty-based state query policy with the same interaction logic as the system-level \textbf{Ours} policy. \textbf{Ours2} used the same state-level query interface with a different WoZ scenario profile to expose participants to a broader range of state-query situations. \textbf{KnowNo} served as an action-level query baseline, in which users were asked to select the robot's next action rather than verify a task-level state hypothesis.

Each domain consisted of five scenarios and five interaction conditions. The assignment between scenarios and conditions was randomized for each participant such that every participant experienced all five conditions exactly once per domain. Detailed scenario descriptions are provided in Appendix~\ref{appx:user}.

\subsubsection{Materials and Procedure}~\label{user:materials_procedure}
Participants interacted with the robot through a tablet-based interface. The interface displayed the current task state, visual sensing information, task-relevant objects, and robot-generated questions with selectable answer options. The waste sorting scenarios required participants to identify waste categories, including general waste, paper, cans, and plastic. The tomato harvesting scenarios required participants to identify tomato states and support harvesting decisions, such as whether a target tomato was present, ripe, unripe, or rotten.

Before the experiment, participants received an explanation of the two task domains and completed practice interactions with the tablet interface. During the main session, they completed ten trials without assistance. If a participant selected an incorrect answer, the response was recorded as incorrect for evaluation purposes, but the scenario continued using the intended task state so that later interactions were not affected by earlier mistakes. After each scenario, participants completed questionnaire items for workload and situation awareness.

\subsubsection{Measures and Analysis}~\label{user:measures}
Task success measured whether participants provided all required responses correctly within a task. \textbf{No} is omitted from task success because it contains no scorable user responses. Operation time measured active user operation time after subtracting passive scenario video duration; \textbf{No} is also omitted from this measure because it involves no active interaction. Workload was measured using NASA Raw-TLX. We still report \textbf{No} for workload because no feedback is expected to impose the least interaction burden, but this condition does not provide user corrections to the robot. Situation awareness was measured using SAGAT-style questions, grouped into Lv1 perception, Lv2 comprehension, and Lv3 projection.

For condition-level comparisons, we used Cochran's $Q$ test for task success and Friedman tests for operation time and workload. Given the small sample size and one trial per condition in each domain, we use these tests to qualify the descriptive results rather than to make broad claims about population-level effects. The statistical tests are exploratory and are reported without correction for multiple comparisons; pairwise post-hoc tests are not used to claim condition-level superiority.

\subsubsection{Task Success and Operation Time}~\label{user:results_task_time}
Table~\ref{tab:user_task_time} and Fig.~\ref{fig:user_task_time_results} summarize task success and operation time. In waste sorting, the proposed state-query conditions achieved high descriptive success rates ($91.7\%$ for \textbf{Ours1} and $100.0\%$ for \textbf{Ours2}) with shorter operation time than \textbf{All}, while \textbf{KnowNo} showed lower task success. In tomato harvesting, \textbf{KnowNo} achieved the highest descriptive success rate, while the two proposed conditions required less operation time than \textbf{All}.



\begin{table*}[t]
\centering
\caption{Task success and operation time by domain and condition. Values are mean $\pm$ standard error.}
\label{tab:user_task_time}
\small
\setlength{\tabcolsep}{5pt}
\begin{tabular*}{\textwidth}{@{\extracolsep{\fill}}llcccc}
\toprule
\textbf{Domain} & \textbf{Measure} & \textbf{All} & \textbf{Ours1} & \textbf{Ours2} & \textbf{KnowNo} \\
\midrule
Waste & Task success (\%) & $91.7 \pm 8.3$ & $91.7 \pm 8.3$ & $100.0 \pm 0.0$ & $8.3 \pm 8.3$ \\
Waste & Operation time (s) & $30.0 \pm 6.7$ & $11.2 \pm 1.4$ & $10.0 \pm 1.3$ & $29.2 \pm 5.7$ \\
Tomato & Task success (\%) & $66.7 \pm 14.2$ & $75.0 \pm 13.1$ & $75.0 \pm 13.1$ & $100.0 \pm 0.0$ \\
Tomato & Operation time (s) & $35.2 \pm 3.2$ & $21.4 \pm 3.5$ & $17.6 \pm 2.2$ & $25.9 \pm 5.4$ \\
\bottomrule
\end{tabular*}
\end{table*}

\subsubsection{Workload and SAGAT}~\label{user:results_workload_sa}
Fig.~\ref{fig:user_workload_results} and Table~\ref{tab:user_workload_sa} summarize NASA Raw-TLX workload. The workload scores were generally low across conditions. The \textbf{No} condition had the lowest workload in both domains, but this condition also provided no corrective support to the robot and is therefore not directly comparable to feedback-based conditions. Among feedback-based conditions, the proposed state-query conditions showed workload values comparable to or lower than \textbf{All} in several cases, without imposing a large subjective burden.

\begin{figure*}[t]
    \centering
    \begin{subfigure}[t]{0.40\textwidth}
        \centering
        \includegraphics[width=\linewidth]{figures/user_study/waste_nasa_rtlx.png}
        \caption{Waste.}
        \label{fig:user_waste_nasa}
    \end{subfigure}
    \hspace{0.03\textwidth}
    \begin{subfigure}[t]{0.40\textwidth}
        \centering
        \includegraphics[width=\linewidth]{figures/user_study/tomato_nasa_rtlx.png}
        \caption{Tomato.}
        \label{fig:user_tomato_nasa}
    \end{subfigure}
    \caption{NASA Raw-TLX workload. Lower values indicate lower workload. Error bars indicate standard error.}
    \label{fig:user_workload_results}
\end{figure*}

\begin{table}[t]
\centering
\caption{NASA Raw-TLX workload by domain and condition. Values are mean $\pm$ standard error.}
\label{tab:user_workload_sa}
\footnotesize
\setlength{\tabcolsep}{2pt}
\begin{tabular*}{\columnwidth}{@{\extracolsep{\fill}}lccccc}
\toprule
\textbf{Domain} & \textbf{All} & \textbf{No} & \textbf{Ours1} & \textbf{Ours2} & \textbf{KnowNo} \\
\midrule
Waste & $2.06 \pm 0.24$ & $1.65 \pm 0.16$ & $2.26 \pm 0.28$ & $1.89 \pm 0.21$ & $2.54 \pm 0.29$ \\
Tomato & $2.21 \pm 0.24$ & $1.72 \pm 0.19$ & $1.89 \pm 0.20$ & $1.81 \pm 0.15$ & $1.81 \pm 0.17$ \\
\bottomrule
\end{tabular*}
\end{table}

\subsubsection{SAGAT Results}~\label{user:results_sagat_levels}
Table~\ref{tab:user_sagat_levels} reports SAGAT accuracy by level. The results are mixed and do not show a consistent advantage for a single condition: some proposed-query conditions scored well on selected levels, but the pattern varied across domains and SAGAT levels. The \textbf{No} condition also obtained relatively high scores in some cases, which should be interpreted cautiously because this condition involved observation without corrective feedback.

\begin{table*}[t]
\centering
\caption{SAGAT level accuracy by domain and condition. Values are mean $\pm$ standard error. Lv1, Lv2, and Lv3 correspond to perception, comprehension, and projection, respectively.}
\label{tab:user_sagat_levels}
\small
\setlength{\tabcolsep}{4pt}
\begin{tabular*}{\textwidth}{@{\extracolsep{\fill}}llccc}
\toprule
\textbf{Domain} & \textbf{Condition} & \textbf{Lv1} & \textbf{Lv2} & \textbf{Lv3} \\
\midrule
Waste & All & $0.92 \pm 0.06$ & $0.38 \pm 0.13$ & $0.38 \pm 0.07$ \\
Waste & No & $0.96 \pm 0.04$ & $0.25 \pm 0.12$ & $0.79 \pm 0.10$ \\
Waste & Ours1 & $0.79 \pm 0.10$ & $0.63 \pm 0.07$ & $0.75 \pm 0.08$ \\
Waste & Ours2 & $0.88 \pm 0.07$ & $0.08 \pm 0.06$ & $0.63 \pm 0.09$ \\
Waste & KnowNo & $0.75 \pm 0.08$ & $0.46 \pm 0.13$ & $0.58 \pm 0.12$ \\
\midrule
Tomato & All & $0.88 \pm 0.07$ & $0.46 \pm 0.10$ & $0.71 \pm 0.10$ \\
Tomato & No & $0.75 \pm 0.08$ & $0.29 \pm 0.10$ & $0.83 \pm 0.09$ \\
Tomato & Ours1 & $0.63 \pm 0.07$ & $0.42 \pm 0.08$ & $0.83 \pm 0.07$ \\
Tomato & Ours2 & $0.67 \pm 0.07$ & $0.46 \pm 0.13$ & $0.96 \pm 0.04$ \\
Tomato & KnowNo & $0.63 \pm 0.07$ & $0.50 \pm 0.09$ & $0.63 \pm 0.07$ \\
\bottomrule
\end{tabular*}
\end{table*}

\subsubsection{Summary}~\label{user:discussion}
This study served as a preliminary usability check for the proposed query interface alongside the main system-level evaluation in Section~\ref{exp:sys_eval}. Participants were able to use state-level queries in both task domains, and the proposed conditions often reduced operation time or workload relative to continuous feedback, though outcomes varied by domain and measure.

The results from the two domains show that state-level queries can be applied to real user interaction without imposing excessive burden on participants. In waste sorting, the proposed conditions achieved task success comparable to continuous feedback while reducing operation time. In tomato harvesting, the action-level baseline \textbf{KnowNo} achieved a high success rate, but the proposed conditions also showed shorter operation time and lower workload, indicating reduced user burden. Thus, the user study provides supplementary evidence that the selective state-query strategy evaluated at the system level can be implemented as a usable human interface without causing excessive interaction burden.
